\documentclass{jsarticle}
\usepackage{ascmac,amsmath,amssymb,amsthm,bm,physics,arydshln,mathcomp,wrapfig}
\usepackage{geometry}
\usepackage[inline]{enumitem} %enumerate環境で横に箇条書きをかけるようにする
\usepackage{ulem} %波線や破線、二重線などを引けるようにする
\usepackage{multirow} %表のところで使う
\geometry{a4paper, left=10mm, right=10mm, top=10mm, bottom=30mm}

\begin{document}

\begin{tabular}{p{1cm}p{2cm}p{0.5cm}p{1.5cm}p{1cm}p{1.5cm}p{0.5cm}p{1cm}p{1cm}p{0.5cm}p{2cm}p{0.1cm}p{1cm}p{1cm}}
        \cline{1-11} \cline{13-14}
\multicolumn{1}{|c|}{\multirow{2}{*}{科目名}} &  \multicolumn{1}{c|}{\multirow{2}{*}{ 解析学2 }} & \multicolumn{1}{c|}{\multirow{2}{*}{対象}} & \multicolumn{1}{c|}{\multirow{2}{*}{ \ 3MA \ }} & \multicolumn{1}{c|}{\multirow{2}{*}{学部}} & \multicolumn{1}{c|}{\multirow{2}{*}{創域理工学部}} & \multicolumn{1}{c|}{\multirow{2}{*}{学科}} & \multicolumn{2}{c|}{\multirow{2}{*}{数理科学科}} & \multicolumn{1}{c|}{学籍} & \multicolumn{1}{c|}{ \ \ \ \ \ \ \ \ \ \ \ \ \ \ \ \ \ \ \ \ \ } & \multicolumn{1}{c|}{} & \multicolumn{1}{c}{ \ \ 評点} & \multicolumn{1}{c|}{}  \\ \cline{13-14}
\multicolumn{1}{|c|}{}  & \multicolumn{1}{c|}{} & \multicolumn{1}{c|}{}  & \multicolumn{1}{c|}{}  & \multicolumn{1}{c|}{} & \multicolumn{1}{c|}{}  & \multicolumn{1}{c|}{}  & \multicolumn{2}{c|}{}  & \multicolumn{1}{c|}{番号}  & \multicolumn{1}{c|}{} & \multicolumn{1}{c|}{} & \multicolumn{1}{c}{} & \multicolumn{1}{c|}{} \\ \cline{1-11}
\multicolumn{4}{|l|}{\multirow{2}{*}{2024年1月10日（金）}} & \multicolumn{1}{c|}{\multirow{2}{*}{担当}} & \multicolumn{1}{c|}{\multirow{2}{*}{立川 篤}} & \multicolumn{1}{c|}{\multirow{2}{*}{学年}} & \multicolumn{1}{c|}{\multirow{2}{*}{ \ \ 年}} & \multicolumn{1}{c|}{\multirow{2}{*}{氏名}} & \multicolumn{2}{c|}{} & \multicolumn{1}{c|}{} & \multicolumn{1}{c}{} & \multicolumn{1}{c|}{} \\ \cline{13-14} 
\multicolumn{4}{|r|}{3回目} & \multicolumn{1}{c|}{} & \multicolumn{1}{c|}{} & \multicolumn{1}{c|}{} & \multicolumn{1}{c|}{} & \multicolumn{1}{c|}{} & \multicolumn{2}{c|}{} \\ \cline{1-13}
\multicolumn{1}{|c|}{} & \multicolumn{1}{c|}{} & \multicolumn{2}{c|}{\multirow{1.8}{*}{持込に関する}} & \multicolumn{9}{c|}{} \\ 
\multicolumn{1}{|c|}{試験} & \multicolumn{1}{c|}{\multirow{2}{*}{{\LARGE60}分}} & \multicolumn{2}{c|}{\multirow{2.5}{*}{注意事項}} & \multicolumn{9}{c|}{筆記用具以外持込不可} \\ 
\multicolumn{1}{|c|}{時間} & \multicolumn{1}{c|}{} & \multicolumn{2}{c|}{} & \multicolumn{9}{c|}{} \\ \cline{3-13}
\multicolumn{1}{|c|}{} & \multicolumn{1}{c|}{} & \multicolumn{2}{c|}{試験時注意事項} & \multicolumn{9}{c|}{} \\
\cline{1-13}
\end{tabular}

注: 以下，扱っている線形空間はすべて実線形空間としてよい（もちろん，しなくてもよい）．

\begin{enumerate}
\item 
Hilbert空間$X, Y$ の積空間$X \times Y$ の元$(x, y)$ に対して
$$
\left(\left(x_{1}, y_{1}\right),\left(x_{2}, y_{2}\right)\right)_{X \times Y}:=\left(x_{1}, x_{2}\right)_{X}+\left(y_{1}, y_{2}\right)_{Y}
$$
とおく. ただし，$\left(x_{1}, x_{2}\right)_{X},\left(y_{1}, y_{2}\right)_{Y}$ は，それぞれ $X, Y$ での内積を表す.

(a) \ $(\cdot, \cdot)_{X \times Y}$ が内積となることを示せ.

(b) \ $X \times Y$ が$(\cdot, \cdot)_{X \times Y}$ に関してヒルベルト空間となることを示せ.

\vspace{13zw}

\item 
$H$ を Hilbert 空間，$\left\{x_{k}\right\}_{k \in \mathbb{N}}$ を $H$ の正規直交系とする．このとき，任意の $x \in H$ に対して，

$$
\sum_{k=1}^{\infty}\left|\left(x, x_{k}\right)\right|^{2} \leq\|x\|^{2}
$$

が成り立つことを示せ．ただし，$(\cdot, \cdot),\|\cdot\|$ は $H$ の内積とその内積から定まるノルムとする. 

\vspace{13zw}

\item
$X, Y$ を Banach 空間とする. $T$ を $X$ から $Y$ への有界線形作用素で，$D(T)=X$ であるとする（ $D(T)$ は $T$ の定義域）.

$$
\mathcal{N}(T)=\{x \in X \ ; \ T x=0\}
$$

とおくと， $\mathcal{N}(T)$ は $X$ の閉部分空間となることを示せ．


\newpage


\item
$\Omega$ を $\mathbb{R}^{n}$ の有界開集合とし，

$$
L=\left\{u(x) \in L^{2}(\Omega) ; \int_{\Omega} u(x) d x=0\right\}
$$

とおく．このとき次の問いに答えよ．\\
（a）$L$ が $L^{2}(\Omega)$ の閉部分空間となることを示せ. \hspace{8zw}
（b）$L$ への射影 $\operatorname{Proj}_{L}$ を求めよ. \ (証明不要) \\
（c）$L^{2}(\Omega)$ において $L^{\perp}$ を求めよ. \ (証明不要)

\vspace{16zw}

\item
$X, Y$ をBanach空間とする.

(a) 線形作用素 $S: D(S) \subset X \to Y$ ($D(S)$ は$S$ の定義域を表す) が閉作用素であることの定義を述べよ.

(b) 有界線形作用素$T:X \to Y$ は閉作用素となることを示せ.

\vspace{18zw}

\item
$X$ をBanach空間とする. このとき, \ 次の問いに答えよ.

(a) $X$ の共役空間$X^*$ の定義を述べよ. \hspace{8zw}
(b) $X$ における弱収束の定義を述べよ.

(c) $X$ における点列$\{x_n\}$ が$\bar{x} \in X$ に強収束していれば, \ $\{x_n\}$ は$\bar{x}$ に弱収束することを示せ.
 
\end{enumerate}


\end{document}
